\clearpage\thispagestyle{empty}
\pdfbookmark[0]{Quellen zu Polarstern und Sonne (18. Juli 1817)}{polaris1817jul-sources}
\begin{center}\textbf{Quellen zu Polarstern und Sonne (18. Juli 1817)}\end{center}
\noindent Diese KI-Übertragung ins heutige Deutsch folgt dem Brief von Carl Friedrich Gauß vom 18. Juli 1817 in den \href{https://gdz.sub.uni-goettingen.de/id/PPN236007467}{\textit{Werken}, Band~VI} (Göttingen, 1874), S.~393--399. Der Erstabdruck erschien als Beitrag XII, als Auszug aus einem Brief an den Direktor der Seeberger Sternwarte, in der von Bernhard von Lindenau und Johann Gottlieb Friedrich Bohnenberger herausgegebenen \textit{Zeitschrift für Astronomie und verwandte Wissenschaften}, Band~IV, Juli und August 1817, S.~\href{https://api.digitale-sammlungen.de/iiif/image/v2/bsb10061006_00119/full/1600,/0/default.jpg}{119}--\href{https://api.digitale-sammlungen.de/iiif/image/v2/bsb10061006_00131/full/1600,/0/default.jpg}{131}. Die Links zum Erstabdruck führen zum Exemplar der Bayerischen Staatsbibliothek.
\par\smallskip\noindent Auf S.~\href{https://api.digitale-sammlungen.de/iiif/image/v2/bsb10061006_00119/full/1600,/0/default.jpg}{119} bezeichnet der Erstabdruck den östlichen Versatz als $1''{,}90$, die Werkausgabe als $1^{\mathrm s}90$. In der Vermutung zum Schwerpunkt der Sonne stehen auf S.~\href{https://api.digitale-sammlungen.de/iiif/image/v2/bsb10061006_00124/full/1600,/0/default.jpg}{124} $\frac{1}{400}$ und $\frac{1}{1000}$, in der Werkausgabe $\frac{1}{100}$ und $\frac{1}{1000}$. Auf S.~\href{https://api.digitale-sammlungen.de/iiif/image/v2/bsb10061006_00126/full/1600,/0/default.jpg}{126} fehlt die Spalte für Bessels Refraktion, und der Barometerstand vom 8.~April ist mit 7,2 angegeben. In der Werkausgabe steht dort 7~2.
\par\smallskip\noindent Auf S.~\href{https://api.digitale-sammlungen.de/iiif/image/v2/bsb10061006_00128/full/1600,/0/default.jpg}{128} endet die zweite Uhrzeitzeile mit 30,81 wie in der Werkausgabe. Die Reduktion für den 20.~Juni steht auf S.~\href{https://api.digitale-sammlungen.de/iiif/image/v2/bsb10061006_00129/full/1600,/0/default.jpg}{129} als 61~30,15, in der Werkausgabe als 61~30~15. Auf S.~\href{https://api.digitale-sammlungen.de/iiif/image/v2/bsb10061006_00130/full/1600,/0/default.jpg}{130} steht für das auf die Mitte der Sternwarte reduzierte Ergebnis der Sonnenbeobachtungen $51^\circ\,31'\,45''{,}21$. Die Werkausgabe druckt $50^\circ\,31'\,45''21$. Auf S.~\href{https://api.digitale-sammlungen.de/iiif/image/v2/bsb10061006_00131/full/1600,/0/default.jpg}{131} steht \textit{Reduction}, in der Werkausgabe \textit{Refraction}. Die Übertragung folgt der Werkausgabe einschließlich ihrer drei eingeklammerten handschriftlichen Aufzeichnungen. Die Fußnote zu Westphals Koeffizienten steht in beiden Ausgaben. KI-Übertragung und KI-Satz.
\clearpage
